% =====================================================================
%  Chapter 3: Contributions
%  Sources: docs/adr/0001-0008, SOTA_Articles_Report.pdf (gaps G1-G8),
%  docs/report/cowork-script.md section 7 (facts and numbers).
% =====================================================================
\chapter{Contributions}
\label{chap-contribution}

% =====================================================================
\section{Introduction and positioning}
\label{sec-c-intro}
% =====================================================================

The analysis of related work in \cref{sec-synthesis} shows a clear
picture. Recent systems combine LLM agents with an LMS
(ELA Tutor~\cite{lopezgoyez2026ela}), describe interoperable architectures
that were never implemented (IJACSA-Arch~\cite{khaddar2026ijacsa}), or model
a rich learner context without any language model
(CA-MA-DRL~\cite{sharif2026camadrl}). None of them grounds specialised LLM
agents on an explicit, shared model of the learner and of the learner's
context (knowledge, activity, temporal and pedagogical constraints) that
adapts the learning path, the content and the feedback together inside a
real LMS. None of them measures the effect of such a system on learning.
This work aims to fill that gap.

This chapter presents what the framework adds to the state of the art.
\Cref{sec-c-gaps} maps each gap identified in \cref{chap-sota} to a
contribution and gives its status. Sections~\ref{sec-c1} to~\ref{sec-c9}
then present each contribution with the same structure: the
\emph{problem} it addresses, the \emph{proposal}, \emph{how it works},
the \emph{evidence} collected so far and its \emph{limits}.
\Cref{sec-c-compare} compares the framework with the closest related work,
and \cref{sec-c-limits} states the limits of the whole work.
\Cref{fig-contribution-map} gives an overview: each gap on the left is
linked to the contributions that address it, and each contribution is
coloured by its status.

\begin{figure}[htbp]
  \centering
  \reportfig{contribution-map}
  \caption{From literature gaps to contributions. Colours give the
  status of each contribution.}
  \label{fig-contribution-map}
\end{figure}

% =====================================================================
\section{From gaps to contributions}
\label{sec-c-gaps}
% =====================================================================

\Cref{tab-gaps-contributions} links the eight gaps of the literature
(G1 to G8) to the contributions of this work (C1 to C10, plus the
Material Improvement agent and the planned evaluation). The status column
is deliberately strict: a contribution is \emph{implemented} only when it
runs in the prototype, \emph{partial} when a stated part is missing, and
\emph{future work} when nothing of it runs yet.

\begin{table}[htbp]
\centering
\caption{Gaps in the literature and the contributions that address them}
\label{tab-gaps-contributions}
\scriptsize
\begin{tabularx}{\textwidth}{@{}L{1.0cm}YL{3.0cm}L{2.6cm}L{1.3cm}@{}}
\toprule
\textbf{Gap} & \textbf{What is missing in the literature} & \textbf{Our contribution} &
  \textbf{Status} & \textbf{Where} \\
\midrule
G1 & No LLM multi-agent system has an explicit, multi-dimensional context model &
  C1 shared, versioned learner + context model; C2 Context agent &
  Implemented (device and affect not modelled) & \ref{sec-c1}, \ref{sec-c2} \\
G2 & LLM agents are disconnected from knowledge tracing &
  C1; C3 BKT behind the agents & Implemented (offline validation is future work) &
  \ref{sec-c1}, \ref{sec-c3} \\
G3 & Adaptation covers one dimension at a time, rarely inside a real LMS &
  C4 specialised agents on one graph; C10 Moodle integration &
  C4 implemented; C10 partial (LTI 1.3 is future work) & \ref{sec-c4}, \ref{sec-c8} \\
G4 & Tutoring is reactive & C2 Context agent; C5 proactive triggers &
  Implemented (4 rule-based triggers) & \ref{sec-c2} \\
G5 & No evidence of learning effect & Learning-gain evaluation &
  Future work (planned sprint) & \ref{sec-c-limits} \\
G6 & Weak teacher control & C6 visible ``why''; C7 teacher in the loop;
  Material Improvement agent & Implemented & \ref{sec-c4}, \ref{sec-c6}, \ref{sec-c7} \\
G7 & Grounding and reliability across agents are vague or unmeasured &
  C6; C8 grounding and guardrails; Material Improvement agent &
  Implemented (a separate verifier agent is future work) &
  \ref{sec-c5}, \ref{sec-c7} \\
G8 & No French or Arabic support; cloud LLMs on student data; local context ignored &
  C8 (local-model option); C9 multilingual & C8 implemented; C9 partial (Arabic is future work) &
  \ref{sec-c5}, \ref{sec-c9} \\
\bottomrule
\end{tabularx}
\end{table}

% =====================================================================
\section{A shared, versioned learner and context model (C1)}
\label{sec-c1}
% =====================================================================

\paragraph{Problem.} In the systems reviewed, the learner is described by
a profile written in natural language by an LLM
(GenMentor~\cite{wang2025genmentor}), by an interaction history and a
two-variable state (ELA Tutor~\cite{lopezgoyez2026ela}), or by engagement
features (IJACSA-Arch~\cite{khaddar2026ijacsa}). The only rich context
model, in CA-MA-DRL~\cite{sharif2026camadrl}, is not used by language
models and was not implemented (gap G1). When several agents keep their
own picture of the learner, they also disagree with each other, and their
decisions cannot be audited.

\paragraph{Proposal.} The framework keeps \emph{one} learner state per
pair (learner, course), shared by every agent. This \emph{shared learner
state} holds four parts: the knowledge of the learner (the BKT mastery of
each section of the course, see \cref{sec-c3}), the misconceptions
observed in wrong answers, interaction statistics, and the context of the
learner. It is stored in PostgreSQL, versioned, and read and written by
every agent through a single component, the \texttt{StateStore}. No agent
keeps a private copy, and every change is traced.
\Cref{fig-multi-agent} shows the agents around this state, with what each
of them reads and writes.

\begin{figure}[htbp]
  \centering
  \reportfig{multi-agent}
  \caption{Specialised agents around one shared, versioned state.}
  \label{fig-multi-agent}
\end{figure}

\paragraph{How it works.} The state of a learner in a course is a row of
the table \texttt{learner\_course\_state}. Its parts are \texttt{jsonb}
columns, mirrored by Pydantic models on the agent side, and the row
carries an integer \texttt{version}. A turn of an agent follows three
steps:
\begin{enumerate}
  \item the agent loads the state and its version
        (\texttt{load\_turn\_context});
  \item it computes its change as a \emph{pure function} from the old
        state to the new state, without writing anything;
  \item \texttt{commit\_turn} applies the change with an update guarded
        by the version read in step~1
        (\texttt{UPDATE \dots{} WHERE version = <version read>}).
\end{enumerate}
If another agent committed in between, the guarded update matches no row.
The store then reads the state again and applies the same function again,
with a bounded number of retries. This is optimistic concurrency: no
update is lost, and nothing is locked while an LLM is generating text. The
messages of the turn, the new state and a \texttt{context\_snapshot} (the
context the agents saw and produced) are written in one transaction, so
the history of the conversation and the state always agree. Every routing
decision and every agent step is also written to an audit table,
\texttt{agent\_event}, keyed by a \texttt{trace\_id}. LangGraph's own state
is only the working memory of one turn, so there is no second copy of the
conversation.

The \emph{context model} has two levels (\cref{tab-context-dimensions}).
The \emph{course context} is the same for every learner and is stored in
its own versioned table, \texttt{course\_context}. The \emph{learner
context} is a part of the shared learner state, so it is versioned and
committed with the same mechanism as the mastery. Each dimension has an
explicit source, which answers the definition of context as the
information that characterises the situation of an
entity~\cite{dey2001context}. \Cref{fig-context-model} shows the sources
that feed each dimension and the components that use it.

\begin{table}[htbp]
\centering
\caption{Dimensions of the context model}
\label{tab-context-dimensions}
\small
\begin{tabularx}{\textwidth}{@{}L{1.5cm}L{2.1cm}YL{2.8cm}@{}}
\toprule
\textbf{Level} & \textbf{Dimension} & \textbf{Content} & \textbf{Source} \\
\midrule
Course & temporal & deadlines (assignment \texttt{duedate}, quiz \texttt{timeclose}) & Moodle web services \\
Course & pedagogical & prerequisite graph (Moodle restrict-access conditions, else course order); teacher objectives & Moodle, teacher \\
Learner & temporal & last access (the latest of Moodle and our app), days inactive, upcoming deadlines not completed & Moodle + app \\
Learner & activity & completion per activity and section (Moodle, plus in-app views for ``view'' activities), Moodle quiz grades, current section, reading time per section & Moodle + app \\
Learner & pedagogical & readiness per section (prerequisites mastered or completed) & derived by the Context agent \\
Learner & knowledge & BKT mastery per section, misconceptions & Diagnosis agent \\
\bottomrule
\end{tabularx}
\end{table}

\begin{figure}[htbp]
  \centering
  \reportfig{context-model}
  \caption{The explicit context model and its sources.}
  \label{fig-context-model}
\end{figure}

The parts of the state are \texttt{jsonb} on purpose: a new agent can
start writing its part of the learner model without a database
migration. Fields that become stable are then promoted to typed
attributes, and to columns if they need an index.

\paragraph{Evidence.} The agents are covered by 159 automated tests
(\texttt{pytest}), in which LLMs are replaced by fake models, so the tests
never call an LLM. The history of a learner can be replayed from
\texttt{context\_snapshot}, one version per turn. In a
live verification on the demo course, a student who had been inactive in
Moodle opened a lesson in the app; the forced context sync then recorded
the activity as completed, set the days of inactivity to zero, moved the
current section forward and marked the next section as ready, all in one
new version of the learner state.

\paragraph{Limits.} Device and affect are not modelled: they are cheap to
add from the web app but were not needed for the thesis. The
\texttt{jsonb} parts trade some type safety in SQL for the ability to
evolve; validation happens in the Pydantic models. The embedding
dimension is fixed in the schema (1024 for bge-m3), so changing the
embedding model requires a migration and a new ingestion of the course.

% =====================================================================
\section{The Context agent and proactive interventions (C2, C5)}
\label{sec-c2}
% =====================================================================

\paragraph{Problem.} The tutors reviewed are reactive: they wait for the
learner's question (gap G4). ELA Tutor chooses \emph{how} to answer a
question, but it does not act when a deadline approaches or when a
learner stops working~\cite{lopezgoyez2026ela}. Proactive triggers exist
only as rules in an unimplemented design~\cite{khaddar2026ijacsa} or as
part of a conceptual framework~\cite{sharif2026camadrl}.

\paragraph{Proposal.} A dedicated \emph{Context agent} implements the
context lifecycle: it acquires data from Moodle and from the app, models
it into the course and learner contexts, and disseminates it to the other
agents. After each update, explicit rules turn the context into
\emph{interventions}: suggestions shown to the learner with their reason
and a suggested action.

\paragraph{How it works.} Moodle offers no web service for activity logs,
and its core is never modified, so the context is \emph{polled} through
web services, always through the API, which stays the only Moodle client
(\cref{sec-c8}). The sources are the course contents with their
restrictions and dates, the completion status of each learner, the list
of enrolled users with their last access, and the grade items. The
Context agent runs in two cases:
\begin{itemize}
  \item on a schedule: a BullMQ repeatable job every 10 minutes for every
        course with ingested material;
  \item on demand: when a learner opens a course in the web app,
        throttled to one sync per minute, and forced after the first view
        of an activity.
\end{itemize}
The sequence of a sync is detailed in \cref{sec-d-dynamic}
(\cref{fig-seq-context}).

Implementing the agent showed how much of the context Moodle actually
exposes, and these findings shaped the design. The service account needs
the capability \texttt{moodle/course:ignoreavailabilityrestrictions}
(read-only, as a teacher); otherwise Moodle evaluates the restrictions as
the service account and hides restricted activities and their grades.
The assignment and quiz web services return nothing for a user who is not
enrolled, so deadlines are read from the module dates of the course
contents. The raw rule of a restriction is not exposed, only its HTML
sentence, from which the referenced activity is read. Finally, the
\emph{current section} is defined as the furthest section with completed
activity: with ``the first unfinished section'', a prerequisite gap could
never be detected.

\Cref{tab-triggers} lists the four triggers. They are rules written in
code; the LLM writes nothing at this step. Each intervention carries a
de-duplication key (trigger, target and period), so a sync never repeats
a suggestion that the learner has already seen or dismissed, and a status
(new, seen, dismissed, done).

\begin{table}[htbp]
\centering
\caption{Proactive triggers of the Context agent (default thresholds)}
\label{tab-triggers}
\small
\begin{tabularx}{\textwidth}{@{}L{3.4cm}YY@{}}
\toprule
\textbf{Trigger} & \textbf{Condition} & \textbf{Suggested action} \\
\midrule
\texttt{deadline\_soon} & a deadline within 3 days, activity not completed & open the activity; a quiz if the section is weak \\
\texttt{failed\_quiz} & Moodle quiz grade below 50\,\% & a quiz on that section \\
\texttt{inactive} & no access for 7 days & the recommended next section \\
\texttt{prerequisite\_gap} & the learner works on a section whose prerequisites are not met & a quiz on the prerequisite \\
\bottomrule
\end{tabularx}
\end{table}

The other agents read the context through the shared state. The
Recommendation agent puts first a section with a close deadline that is
not mastered, and places an unmet prerequisite before the section that
needs it; a close deadline comes before an unmet prerequisite, because
the due date does not move. The Tutor receives a one-line summary of the
context, for example a deadline in two days on the current section, so
that it can relate its answer to what the learner is doing.

\paragraph{Evidence.} The class view of the teacher reuses the same
trigger rules as risk flags (\cref{sec-c6}). In the verification of the
teacher pages on the demo course, each of the three students showed the
expected flags: close deadline, failed Moodle quiz, prerequisite gap,
inactivity and below objective. In a live tutor turn, the \texttt{adapt}
step reported the deadline of an assignment due in two days on the
current section (\cref{sec-c4}).

\paragraph{Limits.} Polling adds a delay of up to the sync interval and a
Moodle load proportional to the number of learners times the number of
courses. A local Moodle plugin for events is the upgrade path. Deadline
extensions granted to one learner are not read yet: the course deadline
applies. The triggers are rules with fixed defaults; they were checked on
the demo course but not tuned on real classes.

% =====================================================================
\section{Knowledge tracing behind the agents (C3)}
\label{sec-c3}
% =====================================================================

\paragraph{Problem.} No reviewed LLM-based system feeds a calibrated
estimate of mastery to its agents (gap G2). Deep knowledge tracing
appears only in CA-MA-DRL, which uses no LLM~\cite{sharif2026camadrl}.
Without such an estimate, an agent cannot decide what the learner should
study next, nor adapt the depth of an explanation.

\paragraph{Proposal.} The framework applies Bayesian Knowledge Tracing
(BKT)~\cite{corbett1994kt} to each section of the Moodle course. A
\emph{topic} is a Moodle section, with the key \texttt{s<n>}: the teacher
already defines sections, they are stable across re-ingestions, and a jury
can read them. Agents receive mastery numbers, not raw text.

\paragraph{How it works.} Each topic carries $P(L)$, the probability that
the learner has learned it. After a graded answer, the posterior is
computed with the slip probability $S$ and the guess probability $G$.
After a correct answer:
\begin{equation}
  P(L \mid \text{correct}) =
  \frac{P(L)\,(1-S)}{P(L)\,(1-S) + \bigl(1-P(L)\bigr)\,G},
  \label{eq-bkt-correct}
\end{equation}
and after an incorrect answer:
\begin{equation}
  P(L \mid \text{incorrect}) =
  \frac{P(L)\,S}{P(L)\,S + \bigl(1-P(L)\bigr)\,(1-G)}.
  \label{eq-bkt-incorrect}
\end{equation}
The learner may then learn the topic at this opportunity, with
probability $T$:
\begin{equation}
  P(L') = P(L \mid \text{obs}) + \bigl(1 - P(L \mid \text{obs})\bigr)\,T.
  \label{eq-bkt-learn}
\end{equation}
\Cref{tab-bkt-params} gives the parameters. A topic counts as mastered
when $P(L) \geq 0.95$, or at the target set by the teacher for that
section (\cref{sec-c6}). Only graded answers move $P(L)$; chat turns are
weak and ambiguous evidence, so they only update the interaction
statistics. A wrong answer also records the misconception attached to the
chosen distractor. The update is a pure function committed through the
\texttt{StateStore}, in the same transaction as the response, and traced
with the values before and after.

\begin{table}[htbp]
\centering
\caption{BKT parameters (global defaults)}
\label{tab-bkt-params}
\small
\begin{tabular}{@{}llc@{}}
\toprule
\textbf{Parameter} & \textbf{Meaning} & \textbf{Value} \\
\midrule
$P(L_0)$ & prior probability that the topic is known & 0.20 \\
$T$ & probability of learning at each opportunity & 0.15 \\
$S$ & slip: wrong answer although the topic is known & 0.10 \\
$G$ & guess: right answer although unknown (4-option item) & 0.25 \\
\bottomrule
\end{tabular}
\end{table}

The Recommendation agent ranks the sections that have ingested material
with explicit rules: sections never assessed first (in course order), then
the lowest $P(L)$ below the mastery threshold, then sections with
recurring misconceptions. Key sections chosen by the teacher come first
among the sections never assessed and among those with low mastery, and
the context rules of \cref{sec-c2} apply. Every
recommendation carries a reason code, translated by the interface.
\Cref{fig-adaptive-loop} shows the resulting adaptive loop: a quiz
produces graded answers, the Diagnosis agent updates the mastery, the
Recommendation agent proposes the next section, and the Context agent
feeds interventions into the same state. The Tutor also reads the mastery
to adapt the depth of its explanations.

\begin{figure}[htbp]
  \centering
  \reportfig{adaptive-loop}
  \caption{The adaptive loop.}
  \label{fig-adaptive-loop}
\end{figure}

\paragraph{Evidence.} The behaviour of BKT with these parameters was
measured by simulation: 2000 simulated learners per profile, 20 answers
each (\cref{tab-bkt-sim}). Learners who know the topic reached mastery
after a median of 3 answers; learners who learn it as the model assumes,
after 7. Learners who only guess were wrongly rated mastered 1.6\,\% of
the time after 5 answers, 2.5\,\% after 10 and 2.4\,\% after 20, and
dropped back below the threshold on their next wrong answer, so the
section was recommended again. Over 20 answers, 21.1\,\% of them crossed
the threshold at least once, after a median of 11 answers: the learning
term $T$ slowly lifts the estimate. A rule ``at least $N$ answers before
mastery'' was considered and rejected: from the prior, three correct
answers already give $P(L) \approx 0.958$ by \crefrange{eq-bkt-correct}{eq-bkt-learn}.

\begin{table}[htbp]
\centering
\caption{BKT simulation (2000 learners per profile, 20 answers each)}
\label{tab-bkt-sim}
\small
\begin{tabular}{@{}lccccc@{}}
\toprule
\textbf{Profile} & \textbf{$P(L)$ after 5} & \textbf{after 10} & \textbf{after 20} &
  \textbf{reached mastery} & \textbf{median answers} \\
\midrule
knows   & 0.96 & 1.00 & 1.00 & 100\,\%  & 3 \\
learns  & 0.57 & 0.81 & 0.96 & 94.5\,\% & 7 \\
guesses & 0.33 & 0.35 & 0.35 & 21.1\,\% (at least once) & 11 \\
\bottomrule
\end{tabular}
\end{table}

\paragraph{Limits.} The parameters are global, not fitted per section.
Fitting $T$ per topic from real answers (for example with
expectation-maximisation) is the principled fix for the guessing profile,
once enough data exists. The offline validation of the model on a public
dataset such as ASSISTments, with the AUC as metric, is future work.

% =====================================================================
\section{Specialised agents on one explicit graph (C4, C6)}
\label{sec-c4}
% =====================================================================

\paragraph{Problem.} Each reviewed system adapts one dimension: ELA Tutor
adapts how to answer, GenMentor adapts the path and the content outside
any LMS, IJACSA-Arch adapts interventions by rules (gap G3). In addition,
when an LLM takes pedagogical decisions, the decisions are hard to
explain to a teacher (gaps G6 and G7).

\paragraph{Proposal.} The framework splits the work among specialised
agents, coordinated through explicit graphs and the shared state
(\cref{tab-agents-roster}). It follows a principle already stated by ELA
Tutor~\cite{lopezgoyez2026ela}: \emph{code decides, the LLM writes}.
Routing, grading, mastery, ranking and triggers are computed in code; the
LLM only classifies when the rules are not confident, and writes text.
Every decision carries a short ``why'' visible to the learner and to the
teacher. There is no separate Planner or Content Generator agent: the
Recommendation agent plans the path, and the Assessment agent generates
the exercises.

\begin{table}[htbp]
\centering
\caption{The agents of the framework: role, who decides, use of the LLM}
\label{tab-agents-roster}
\scriptsize
\begin{tabularx}{\textwidth}{@{}L{2.4cm}YL{3.3cm}L{2.6cm}@{}}
\toprule
\textbf{Agent} & \textbf{Role} & \textbf{Decision by} & \textbf{LLM use} \\
\midrule
Orchestrator & classify the message, route the turn & rules first; LLM classifier only when the rules are not confident & structured classification (fallback) \\
Tutor & answer from the course & similarity gate (code); hint mode for exercises & writes the answer \\
Assessment & quizzes & bank first; review mode & writes missing questions (validated in code) \\
Diagnosis & knowledge tracing & BKT (code) & none \\
Recommendation & what to study next & rules (code) & none \\
Context & context model and interventions & rules (code) & none \\
Learning Analytics & class dashboard & code & none \\
Material Improvement & improve the material & 5 detectors (code, statistics) & writes diagnosis and draft \\
Ingestion (worker) & index the course & code & embeddings only \\
\bottomrule
\end{tabularx}
\end{table}

\paragraph{How it works.} A tutor turn is a LangGraph~\cite{langgraph_docs}
graph (\cref{fig-tutor-graph}). The \texttt{classify} node finds the intent
of the message (concept, exercise, quiz, course\_info, smalltalk or
off\_topic), its language and its topic, with rules first. Smalltalk ends
the turn. Otherwise \texttt{retrieve} searches the course
(\cref{sec-c5}), and a gate written in code chooses the next node: a quiz
intent leads to \texttt{suggest\_quiz}; a best similarity at or above the
threshold leads to \texttt{adapt} then \texttt{generate}; a lower
similarity leads to \texttt{not\_found}. The \texttt{adapt} step maps the
mastery of the section to a level: below 0.4 the explanation is low level
(step by step), below 0.8 medium, otherwise high; the level is
\emph{unknown} if the section was never assessed.

\begin{figure}[htbp]
  \centering
  \reportfig{tutor-graph}
  \caption{Graph of a tutor turn: code decides, the LLM writes.}
  \label{fig-tutor-graph}
\end{figure}

Each node emits an \emph{agent step}, first \emph{running} then
\emph{done} with its details and its duration, streamed to the chat as a
server-sent event. The details are the values that the node already
records in its \texttt{agent\_event} span, so what the student sees is
what the trace holds. The steps are stored with the answer, so a reopened
conversation shows them again. The model's own chain of thought is
\emph{never} shown: in hint mode it could reveal the solution that the
tutor must withhold, and it is not a trace of the system's decisions.

\paragraph{Evidence.} For the question ``What does \texttt{range(5)} give
in a \texttt{for} loop?'', the chat showed: \texttt{classify} (concept,
by rule, 7\,ms), \texttt{retrieve} (6 passages, best section ``Conditions
and loops'' with similarity 0.62, 3.4\,s), \texttt{adapt} (mastery 20\,\%,
step by step, deadline in 2 days) and \texttt{generate} (1.2\,s, with a
citation). An off-topic question about Django and PostgreSQL ended in
\texttt{not\_found}, with the closest sections of the course. The steps
make the multi-agent decisions visible at no extra LLM cost.

\paragraph{Limits.} The rules of the classifier and the thresholds were
set on one demo course. The LLM classifier is only a fallback, and the
router is the extension point for a learned policy such as the
reinforcement-learning meta-agent of ELA Tutor~\cite{lopezgoyez2026ela},
which is out of scope here.

% =====================================================================
\section{Grounding and guardrails (C8)}
\label{sec-c5}
% =====================================================================

\paragraph{Problem.} Grounding on the course material is vague in the
reviewed systems (``curated sources'' in ELA Tutor) or only described
conceptually (IJACSA-Arch), and the cascade of hallucinations from one
agent to the next, which the survey of LLM-based multi-agent systems warns
about~\cite{guo2024mas}, is not measured (gap G7). Most systems also send
student data to proprietary cloud models (gap G8), which matters under
the Algerian law on personal data~\cite{algeria_law1807}.

\paragraph{Proposal.} The framework grounds every generated text in the
course and checks it in code before it reaches a learner. A local model
can replace the hosted one.

\paragraph{How it works.}
\begin{itemize}
  \item \emph{Hybrid retrieval.} The query is matched by vector similarity
        (pgvector, cosine) and by PostgreSQL full-text search, and the two
        rankings are fused with reciprocal rank fusion~\cite{cormack2009rrf}.
        The 6 best passages are kept. The embeddings are computed with the
        multilingual model bge-m3~\cite{chen2024bgem3}.
  \item \emph{Similarity gate.} An answer is generated only when the best
        similarity reaches the threshold (0.42 by default, configurable).
  \item \emph{Citations checked in code.} The answer cites passages by
        number; the numbers are validated against the retrieved passages.
  \item \emph{Refusal.} Below the threshold, the tutor says that the
        answer is not in the course and lists the closest sections.
  \item \emph{Hint mode.} For an exercise, the tutor gives hints, not the
        solution.
  \item \emph{Validated generation.} A generated quiz item is kept only if
        it has four distinct options, a valid answer index, sources within
        the retrieved chunks and a unique stem. A draft of the Material
        Improvement agent that quotes a passage keeps the quote only if it
        appears in the excerpts.
  \item \emph{Local model.} The chat model can run locally with Ollama
        (\texttt{qwen3:4b-instruct}), so student data need not leave the
        institution.
\end{itemize}

\paragraph{Evidence.} The off-topic question of \cref{sec-c4} was refused
with the closest sections. For quiz generation on the four sections of
the demo course (writer \texttt{openai/gpt-oss-20b} on Groq), the four
quizzes were generated on the first attempt, 18 of the 20 drafted items
(90\,\%) passed the validation in code, the median time was 3.1\,s per
quiz, and the positions of the correct option were spread over the four
letters (A/B/C/D = 4/3/5/6). Every distractor carried a misconception label. The
mean Ragas faithfulness~\cite{es2024ragas} of the explanations was 0.54.
This value mixes real factual slips with correct general knowledge that
is not in the excerpts, and with false negatives when the material and
the questions are in different languages. Faithfulness alone therefore
cannot certify an item, which is why generated questions go through a
teacher review (\cref{sec-c6}).

\paragraph{Limits.} There is no separate verifier agent that checks every
answer before delivery; it is future work. The rate of hallucinations
cascading between agents is not measured yet.

% =====================================================================
\section{The teacher in the loop (C7)}
\label{sec-c6}
% =====================================================================

\paragraph{Problem.} Teacher control is future work in ELA Tutor, limited
to dashboards in IJACSA-Arch, absent in GenMentor and unimplemented in
CA-MA-DRL (gap G6). The survey of LLMs for education calls for the
instructor to act as a special agent of the system~\cite{wang2024llm4edu}.
The faithfulness measured in \cref{sec-c5} also shows that generated
questions must be checked by a human before they count.

\paragraph{Proposal.} The teacher controls three things: a reviewed
\emph{question bank}, \emph{objectives} for each section, and a
\emph{class dashboard}.

\paragraph{How it works.}
\begin{itemize}
  \item \emph{Question bank.} A quiz is filled from the bank first:
        approved questions of the section, in the learner's language, that
        this learner has never answered. Only the shortfall is generated
        by the LLM. A per-course \emph{review mode} decides what happens to
        generated questions: in \texttt{review} mode (the default) they
        enter the bank as \emph{pending}, and the quiz is served only if
        enough questions are approved; in \texttt{auto\_publish} mode they
        are approved at once and used. The teacher can generate, edit,
        approve or reject questions; an edit goes through the same checks
        as an LLM draft, and every review is traced. A quiz item is a copy
        of a bank question, so a later edit never changes a quiz already
        taken or its BKT history.
  \item \emph{Objectives.} For each section, the teacher sets a target
        mastery and marks key sections. The Recommendation agent, the
        readiness computed by the Context agent and the dashboard all read
        these objectives; none of them copies them.
  \item \emph{Class dashboard.} The Learning Analytics agent computes, from
        the shared state of every enrolled student, a view per section
        (learners assessed, mastered, mean mastery) and per student
        (mastery, inactivity, current section, open suggestions and risk
        flags). The risk flags reuse the trigger rules of
        \cref{tab-triggers}, so ``at risk'' has one definition, plus a flag
        for a key section below its target.
\end{itemize}
The life cycles of a bank question and of a material suggestion are shown
in \cref{fig-states} (\cref{chap-implementation}).

\paragraph{Evidence.} The flow was verified through the public API on the
demo course, starting from an empty bank in \texttt{review} mode. A
student's quiz was first refused with the code \texttt{awaiting\_review},
while questions were written in the background for the teacher. The
teacher edited one question, rejected one and approved three; an edit
with duplicate options was refused. The student's next quiz was served
from the approved questions without any LLM call, and never contained the
rejected question. One of the generated questions claimed that
\texttt{if}/\texttt{elif}/\texttt{else} branches are tested ``from the
most restrictive to the least'', which the course does not say: exactly
the kind of item the review exists to catch.

\paragraph{Limits.} In \texttt{review} mode a learner may have to wait
for the teacher; the teacher can fill the bank in advance. Rejecting a
question that was already answered does not replay the BKT history.

% =====================================================================
\section{An evidence-based Material Improvement agent}
\label{sec-c7}
% =====================================================================

\paragraph{Problem.} The systems reviewed help the learner, but none uses
the learners' difficulties to point the teacher to the parts of the
course that cause them. An LLM asked to ``critique the course'' would
produce plausible but unverifiable advice (gaps G6 and G7).

\paragraph{Proposal.} The Material Improvement agent goes beyond the
contributions planned in the literature analysis. It applies the same
principle as the other agents: \emph{code finds the problems with
statistics, the LLM writes a draft, the teacher decides}. The app never
writes to Moodle: the agent only suggests. It runs on demand, when the
teacher clicks a button, because a scheduled run would waste the quota of
a free-tier LLM.

\paragraph{How it works.} \Cref{fig-material-pipeline} shows the
pipeline. The evidence is loaded from the shared state: answers with the
misconception label of each wrong option, the tutor's turns that the
material could not answer, the reading time and mastery per section, and
the activity views. Five detectors, written as pure functions of this
evidence, raise issues; each issue carries its numbers, a severity and a
stable key, so a new run never re-creates a rejected suggestion.

\begin{figure}[htbp]
  \centering
  \reportfig{material-pipeline}
  \caption{The Material Improvement agent.}
  \label{fig-material-pipeline}
\end{figure}

\begin{itemize}
  \item \texttt{misconception}: a misconception chosen by at least 30\,\%
        of the learners who answered the section, and more often than
        chance among the wrong answers to the questions that offer it
        (one-sided binomial test, $\alpha = 0.01$, Bonferroni correction
        over the labels of the section);
  \item \texttt{problem\_question}: a bank question with at least
        8 answers and at most 25\,\% correct; a wrong option with
        \emph{positive} discrimination against the rest score (chosen by
        the learners who master the section) points to a wrong key;
  \item \texttt{unanswered}: tutor questions the material could not
        answer, grouped by meaning (cosine $\geq 0.70$ to the centroid of
        the group) and asked by at least 2 learners;
  \item \texttt{unclear\_section}: at least 40\,\% of the assessed
        learners read the section for 5 minutes or more and stay below
        0.5 mastery;
  \item \texttt{skipped\_activity}: an activity opened by less than
        30\,\% of the learners assessed on its section.
\end{itemize}
For each new issue, the LLM writes, from the evidence and the excerpts of
the section, a diagnosis and a proposed change in the teacher's language,
and a draft (a passage, an example or a corrected question) in the
language of the material. A quoted passage that is not in the excerpts is
dropped. Each suggestion is then \emph{pending}, \emph{accepted} or
\emph{rejected}, and \emph{applied} once the teacher has changed the
material in Moodle; every review is traced.

\paragraph{Evidence.} The detectors were measured on simulated classes
with planted flaws: 200 random classes for each size, with a confusing
section, a content gap, a skipped activity and a question with a wrong
key (\cref{tab-material-eval}). Precision is the share of raised issues
that match a planted flaw.

\begin{table}[htbp]
\centering
\caption{Material Improvement detectors on simulated classes (200 classes per size)}
\label{tab-material-eval}
\small
\begin{tabular}{@{}lccc@{}}
\toprule
& \textbf{10 learners} & \textbf{20 learners} & \textbf{40 learners} \\
\midrule
confusing section found & 70\,\% & 93\,\% & 100\,\% \\
content gap found & 90\,\% & 100\,\% & 100\,\% \\
skipped activity found & 94\,\% & 99\,\% & 100\,\% \\
wrong answer key found and diagnosed & 40\,\% & 86\,\% & 96\,\% \\
precision of raised issues & 89\,\% & 97\,\% & 100\,\% \\
clean section wrongly flagged & 4\,\% & 2\,\% & 0\,\% \\
\bottomrule
\end{tabular}
\end{table}

These results come from a second version of the detectors. The first
version reached a precision of only 32\,\% on classes of 10 learners and
56\,\% on classes of 20: in a small class, three learners making the same
mistake happens by chance. Three rules were changed after this
measurement. The \texttt{misconception} detector now also requires the
binomial test. The \texttt{problem\_question} detector no longer trusts
one dominant wrong option, which is also what a strong misconception
looks like; it uses item discrimination instead, which works even when
nobody chose the key. The grouping threshold of unanswered questions was
calibrated on bge-m3: each paraphrase of the planted gap is at least 0.74
from the centroid of the others, and unrelated questions are at most 0.66
apart. The detectors are now conservative on small classes, as intended:
they miss more than they invent.

The agent was then run live on the demo course with 20 simulated learners
(426 answers through the Assessment agent). All four planted flaws were
found in 1.5\,s, with no false positive. For the wrong key, only 6\,\% of
the answers were correct, the discrimination of the key was $-0.20$, and
the probable right answer was identified. The misconception ``\texttt{range(n)}
includes \texttt{n}'' was shared by 15 of 19 learners ($p = 6 \cdot 10^{-8}$).
The live run also found defects that the tests with fake models could not
show, such as diagnoses written in the language of the excerpts instead
of the teacher's language; they were fixed.

\paragraph{Limits.} The thresholds are defaults for classes of 20 to 40
learners; small classes produce no suggestion rather than noisy ones.
Applying a change remains manual (in Moodle, followed by a new
ingestion), and whether a change helped (mastery before and after) is left
to the evaluation sprint.

% =====================================================================
\section{Moodle integration without modifying its core (C10)}
\label{sec-c8}
% =====================================================================

\paragraph{Problem.} The closest system integrates Moodle through a
low-code orchestration engine~\cite{lopezgoyez2026ela}; the most complete
interoperable design, based on LTI 1.3 and xAPI, was not
implemented~\cite{khaddar2026ijacsa} (gap G3). An institution will not
adopt a framework that modifies the core of its LMS.

\paragraph{Proposal.} Moodle is the LMS backend and is never modified.
The API gateway is its only client and holds the only token. The
integration is read-only: the framework never writes to Moodle. This
contribution is \emph{partial}: the web services and the in-app reader are
implemented, LTI 1.3 single sign-on is future work.

\paragraph{How it works.} \Cref{fig-moodle-integration} shows the
integration. The API reads Moodle through its REST web services (the
functions are listed in the figure). Students read the course inside the
app: a page is fetched with the service token and \emph{sanitised} in the
API (an allow-list of tags and attributes, no scripts, styles or
iframes); files are served through \emph{signed links}, an HMAC bound to
the user, the file path and an expiry of one hour, because a browser
cannot send a token for an image or an embedded PDF. Moodle's
restrict-access conditions that require completing an activity are
honoured for students; other conditions (dates, grades) are shown but not
enforced, and teachers are never locked out.
Views in the app are recorded in our database and counted as completion
for the activities whose only Moodle completion rule is ``view''; this
keeps the context of \cref{sec-c2} correct for a learner who never opens
Moodle.

\begin{figure}[htbp]
  \centering
  \reportfig{moodle-integration}
  \caption{Moodle as a read-only backend.}
  \label{fig-moodle-integration}
\end{figure}

Writing completion back to Moodle was considered and rejected: the only
web service for it is a teacher override, recorded as such in Moodle's
reports, and it would break the read-only principle.

\paragraph{Evidence.} In a live check, a lesson was served as sanitised
HTML, a PDF of the course was streamed through a signed link, an activity
locked by a restriction was returned as locked with Moodle's sentence, and
after the lesson was viewed in the app, the next section was unlocked.
Moodle itself still showed the lesson as not viewed, as intended. The API
tests (80, with Vitest) replay recorded Moodle responses.

\paragraph{Limits.} Without LTI 1.3 there is no single sign-on: the
app's login is separate from Moodle's. Moodle's own quizzes and
assignments still open in Moodle. Moodle's completion report does not show
the views made in the app; our dashboard does.

% =====================================================================
\section{Multilingual support (C9)}
\label{sec-c9}
% =====================================================================

\paragraph{Problem.} No reviewed system supports French or Arabic
(gap G8), whereas Algerian universities teach in Arabic, French and
English.

\paragraph{Proposal and how it works.} The interface is bilingual
(French and English, with react-i18next). The Tutor answers in the
learner's language, quizzes are generated in the learner's language, and
the course material can be in French or in English: the language of each
document is detected at ingestion, and bge-m3 is a multilingual embedding
model, so a question in one language can be compared with a passage in the
other.
The Material Improvement agent writes its diagnosis in the teacher's
language and its draft in the language of the material.

\paragraph{Status and limits.} This contribution is \emph{partial}:
Arabic is future work. Mixing languages also affects the evaluation: the
faithfulness metric produced false negatives when the material and the
questions were in different languages (\cref{sec-c5}).

% =====================================================================
\section{Comparison with related work}
\label{sec-c-compare}
% =====================================================================

\Cref{tab-related-compare} compares the framework with the four primary
works of \cref{sec-related-work}, on the criteria of the gaps. The
comparison of the other works uses only what their articles report; for
this work, it uses the status of \cref{tab-gaps-contributions}.

\begin{table}[htbp]
\centering
\caption[Comparison with related work]{Comparison with related work
(\cyes: yes, \cpartly: partial or design only, \cno: no, n/r: not reported, n/a: not applicable)}
\label{tab-related-compare}
\scriptsize
\begin{tabularx}{\textwidth}{@{}L{3.3cm}*{5}{>{\centering\arraybackslash}X}@{}}
\toprule
\textbf{Criterion} & \textbf{ELA Tutor} & \textbf{IJACSA-Arch} & \textbf{GenMentor} &
  \textbf{CA-MA-DRL} & \textbf{This work} \\
\midrule
LLM agents & \cyes & \cpartly & \cyes & \cno & \cyes \\
Explicit context model & \cpartly & \cpartly & \cpartly & \cpartly & \cyes \\
Knowledge tracing & \cno & \cpartly & \cno & \cpartly & \cyes \\
Proactive & \cno & \cpartly & \cno & \cpartly & \cyes \\
Teacher control & \cpartly & \cpartly & \cno & \cpartly & \cyes \\
Grounding with citations & \cpartly & \cpartly & n/r & n/a & \cyes \\
Real LMS integration & \cyes & \cpartly & \cno & \cpartly & \cpartly \\
Evaluation of learning & \cno & \cno & \cno & \cno & \cno \\
\bottomrule
\end{tabularx}
\end{table}

\todonote{Check each cell of \cref{tab-related-compare} against Chapter 2
once its related-work section is written. ``Partial'' for CA-MA-DRL means
designed but not implemented; ``partial'' for this work's LMS integration
means LTI 1.3 is pending.}

Two lines of the table deserve a comment. The framework is the only one
that combines an explicit context model with knowledge tracing behind LLM
agents, and the only one with implemented proactive triggers and a
reviewed question bank. On the other hand, like every other work, it has
not yet measured a learning gain. The comparison therefore supports the
claim of the framework on design and implementation, not yet on learning
outcomes.

% =====================================================================
\section{Limits and conclusion}
\label{sec-c-limits}
% =====================================================================

The framework has the following limits, which the defence must state
plainly:
\begin{itemize}
  \item \emph{No measured learning gain} (gap G5 remains open). The
        planned evaluation has three parts: an ablation with simulated
        learners (the full framework, the same agents without the context
        model, and a single LLM chatbot); a pilot with real students using
        a pre-test and a post-test, a comparison group, the normalised
        learning gain~\cite{hake1998gain} and the SUS
        questionnaire~\cite{brooke1996sus}; and technical metrics (offline
        BKT AUC, latency, cost, and the correctness of generated content
        rated by teachers).
  \item \emph{One demo course} and \emph{simulated classes}: the
        thresholds and the detectors were measured on simulated learners,
        not on real classes.
  \item \emph{LTI 1.3} single sign-on and \emph{Arabic} support are
        pending.
  \item \emph{Cost and latency} of hosted LLMs: the free tier of the
        provider shaped several choices (on-demand analysis, bank-first
        quizzes), and retrieval dominates the waiting time because the
        embeddings are computed on the CPU.
\end{itemize}

In summary, this work contributes a shared, versioned model of the
learner and of the learner's context, maintained by a Context agent and
fed by knowledge tracing; specialised agents that decide in code and
write with an LLM, with every decision visible; grounding and guardrails
checked in code; a teacher in the loop for generated questions and
course objectives; an evidence-based Material Improvement agent whose
detectors were measured and corrected; and an integration with Moodle
that never modifies its core. The next chapter presents the design that
implements these contributions.
